% ==================================================================
%  NeighborCheck.tex  --  Section 3.2 header "Server-side Computation"
%                         and 3.2.1 "Neighbor Check"
%  Loaded from main.tex.
%  IEEEtran compsoc conference numbers subsubsections (secnumdepth = 3),
%  so 3.2.1--3.2.4 are numbered without changing secnumdepth.
%  Inside 3.2.x, run-in headings use \paragraph.
% ==================================================================

\subsection{Server-side Computation}
\label{sec:server}

The server receives $\mathsf{ct}_1, \dots, \mathsf{ct}_d$ and the
public parameters, and evaluates DBSCAN in three stages, namely
Neighbor Check (Section~\ref{sec:neighbor}), Core Check
(Section~\ref{sec:core}) and
Label Propagation (Section~\ref{sec:lp}). Section~\ref{sec:correct}
analyzes the accuracy of the comparisons they use.

\subsubsection{Neighbor Check}
\label{sec:neighbor}

Neighbor Check decides, for every pair of points that the ordering
leaves in range, whether their Euclidean distance is at most
$\varepsilon_n$. The distance itself needs a square root, which CKKS can
evaluate only through another approximation. The server therefore
compares squared distances: $\|\mathbf{x}_i - \mathbf{x}_j\| \le
\varepsilon_n$ if and only if $\|\mathbf{x}_i - \mathbf{x}_j\|^2 \le
\varepsilon_n^2$.

The server decides one offset $k$ at a time. Rotating a coordinate
ciphertext by $k$ puts the coordinate of point $i+k$ in slot~$i$.
Subtracting the rotated ciphertext from the original one, squaring, and
summing over the $d$ dimensions gives
\begin{equation}
  \label{eq:dist}
  \mathsf{dist}_k^{2} = \sum_{j=1}^{d}
    \bigl( \mathsf{ct}_j - \mathsf{Rot}_k(\mathsf{ct}_j) \bigr)^{2}.
\end{equation}
This is a single ciphertext. Its slot~$i$ holds
$\|\mathbf{x}_i - \mathbf{x}_{i+k}\|^2$, the squared distance between
point~$i$ and point~$i+k$, for every $i$ at once.

The input of $\mathsf{Sgn}$ has to lie in $[-1,1]$. After the
normalization of Section~\ref{sec:client}, every squared distance lies
in $[0,d]$, so the server divides its difference to the threshold by
\begin{equation}
  \label{eq:bound}
  B = 1.05 \cdot d .
\end{equation}
The factor $1.05$ is chosen heuristically. It keeps the scaled input at
most $1/1.05$ in magnitude,
about $5\%$ inside the ends $\pm 1$ of the interval on which the
approximation is fitted. The minimax polynomials grow quickly outside
that interval, so an input that exceeds $1$ even slightly, for instance
by the noise that $\mathsf{dist}_k^{2}$ carries, can come out with the
wrong sign. Label Propagation (Section~\ref{sec:lp}) divides by $1.1$
instead, also a heuristic choice, and leaves a margin of about $10\%$. There the difference of two
labels, which can reach $\pm 1$, is refreshed by bootstrapping before its
sign is taken, and bootstrapping adds more noise than the computation of
$\mathsf{dist}_k^{2}$, which reaches $\mathsf{Sgn}$ without one. The
scaled input is
\begin{equation}
  \label{eq:cmpinput}
  \mathbf{x}_k = \tfrac{1}{B}
    \bigl( \mathsf{dist}_k^{2} - \theta \bigr),
\end{equation}
where $\theta$ is a threshold slightly above $\varepsilon_n^2$, fixed in
Eq.~\eqref{eq:theta}.

The sign of slot~$i$ of $\mathbf{x}_k$ is negative exactly when points
$i$ and $i+k$ are neighbors. Mapping that sign to $\{0,1\}$ gives the
adjacency vector
\begin{equation}
  \label{eq:adj}
  \mathsf{adj}_k = \tfrac{1}{2}
    \bigl( 1 - \mathsf{Sgn}(\mathbf{x}_k) \bigr),
\end{equation}
whose slot~$i$ records whether points $i$ and $i+k$ are neighbors:
\[
  \mathsf{adj}_k[i] =
  \begin{cases}
    1 & \text{if } \mathsf{dist}_k^{2}[i] \le \varepsilon_n^2
        \quad \text{(neighbors)}, \\[2pt]
    0 & \text{if } \mathsf{dist}_k^{2}[i] > \varepsilon_n^2
        \quad \text{(not neighbors)}.
  \end{cases}
\]
$\mathsf{adj}_k$ is the mask vector that Core Check uses next: adding
the masks of all offsets counts the neighbors of every point.

Algorithm~\ref{alg:neighbor} collects these steps, and the server runs
it for $k = 1, \dots, k_{\max}$. The same pairs seen from their other end
need no further comparison. The vector whose slot~$i$ marks the pair
$(i, i-k)$ is a rotation,
\[
  \mathsf{adj}_{n-k} = \mathsf{Rot}_{n-k}(\mathsf{adj}_{k}),
\]
since slot~$i$ of the right-hand side holds $\mathsf{adj}_k[i-k]$. The
whole stage thus costs $k_{\max}$ calls to $\mathsf{Sgn}$, against the
$n(n-1)/2$ pairwise decisions it stands for.

\begin{algorithm}[t]
\caption{\textsf{Neighbor\_check}}
\label{alg:neighbor}
\begin{algorithmic}[1]
\Require coordinate ciphertexts $\mathsf{ct}_1, \dots, \mathsf{ct}_d$;
  offset $k$; public $\theta$, $d$, $n$
\Ensure $\mathsf{adj}_k$, slot $i$ equal to $1$ iff
  $\|\mathbf{x}_i - \mathbf{x}_{i+k}\| \le \varepsilon_n$
\State $B \gets 1.05 \cdot d$
\State $\mathsf{dist}^{2} \gets
  \bigl( \mathsf{ct}_1 - \mathsf{Rot}_k(\mathsf{ct}_1) \bigr)^{2}$
\For{$j = 2$ \textbf{to} $d$}
  \State $\mathsf{dist}^{2} \gets \mathsf{dist}^{2} +
    \bigl( \mathsf{ct}_j - \mathsf{Rot}_k(\mathsf{ct}_j) \bigr)^{2}$
\EndFor
\State $\mathbf{x} \gets \tfrac{1}{B}
  \bigl( \mathsf{dist}^{2} - \theta \bigr)$
\State $\mathsf{adj}_k \gets
  \tfrac{1}{2}\bigl( 1 - \mathsf{Sgn}(\mathbf{x}) \bigr)$
\State \Return $\mathsf{adj}_k$
\end{algorithmic}
\end{algorithm}
